\documentclass[11pt,a4paper]{article}
\usepackage[top=2.3cm,bottom=2.3cm,left=2.2cm,right=2.2cm]{geometry}

\usepackage{amsmath,amssymb}
\usepackage{fontspec}
\setmainfont{Liberation Serif}
\setsansfont{Liberation Sans}
\setmonofont{DejaVu Sans Mono}[Scale=0.84]

\usepackage{xcolor}
\definecolor{brandblue}{RGB}{20, 55, 105}
\definecolor{accentblue}{RGB}{35, 95, 165}
\definecolor{codebg}{RGB}{248, 249, 251}
\definecolor{codeframe}{RGB}{215, 222, 230}
\definecolor{javakeyword}{RGB}{0, 45, 170}
\definecolor{javastring}{RGB}{5, 120, 20}
\definecolor{javacomment}{RGB}{120, 120, 120}
\definecolor{javanumber}{RGB}{20, 75, 220}

\definecolor{noteborder}{RGB}{30, 110, 65}
\definecolor{notebg}{RGB}{242, 249, 245}
\definecolor{warnborder}{RGB}{185, 75, 10}
\definecolor{warnbg}{RGB}{254, 248, 240}
\definecolor{tipborder}{RGB}{30, 95, 160}
\definecolor{tipbg}{RGB}{242, 247, 254}

\usepackage{fancyhdr}
\usepackage{lastpage}
\pagestyle{fancy}
\fancyhf{}
\lhead{\parbox[t]{0.58\textwidth}{\raggedright\small\sffamily\color{brandblue}\textbf{T20 --- Java I/O, Streams, File, NIO.2 e Serializzazione}}}
\rhead{\small\sffamily\color{gray}Programmazione II -- UniVR (2025/2026)}
\cfoot{\small\sffamily Pagina \thepage\ di \pageref{LastPage}}
\renewcommand{\headrulewidth}{0.5pt}
\renewcommand{\headrule}{\hbox to\headwidth{\color{brandblue!70!black}\leaders\hrule height \headrulewidth\hfill}}

\usepackage{titlesec}
\titleformat{\section}{\Large\bfseries\sffamily\color{brandblue}}{\thesection}{0.8em}{}[\titlerule]
\titleformat{\subsection}{\large\bfseries\sffamily\color{accentblue}}{\thesubsection}{0.8em}{}
\titleformat{\subsubsection}{\normalsize\bfseries\sffamily\color{brandblue!85!black}}{\thesubsubsection}{0.8em}{}

\usepackage{needspace}
\usepackage{fancyvrb}
\DefineVerbatimEnvironment{asciibox}{Verbatim}{
  fontsize=\fontsize{7.5pt}{9.2pt}\selectfont,
  frame=single,
  rulecolor=\color{codeframe},
  fillcolor=\color{codebg},
  framesep=2.5mm
}
\usepackage{listings}
\lstset{
  backgroundcolor=\color{codebg},
  frame=single,
  rulecolor=\color{codeframe},
  basicstyle=\ttfamily\footnotesize,
  keywordstyle=\color{javakeyword}\bfseries,
  stringstyle=\color{javastring},
  commentstyle=\color{javacomment}\itshape,
  numberstyle=\tiny\color{gray},
  numbers=left,
  numbersep=7pt,
  stepnumber=1,
  breaklines=true,
  breakatwhitespace=true,
  showstringspaces=false,
  tabsize=4,
  xleftmargin=12pt,
  framexleftmargin=12pt
}

\newcommand{\passthrough}[1]{#1}
\providecommand{\tightlist}{\setlength{\itemsep}{0pt}\setlength{\parskip}{0pt}}

\usepackage{tcolorbox}
\tcbuselibrary{skins,breakable}
\newtcolorbox{studynote}[1][Nota]{
  enhanced,
  breakable,
  colback=notebg,
  colframe=noteborder,
  fonttitle=\bfseries\sffamily\small,
  coltitle=white,
  title={#1},
  arc=2mm,
  boxrule=0.8pt,
  left=9pt, right=9pt, top=7pt, bottom=7pt
}

\newtcolorbox{studywarning}[1][Attenzione]{
  enhanced,
  breakable,
  colback=warnbg,
  colframe=warnborder,
  fonttitle=\bfseries\sffamily\small,
  coltitle=white,
  title={#1},
  arc=2mm,
  boxrule=0.8pt,
  left=9pt, right=9pt, top=7pt, bottom=7pt
}

\newtcolorbox{studytip}[1][Suggerimento]{
  enhanced,
  breakable,
  colback=tipbg,
  colframe=tipborder,
  fonttitle=\bfseries\sffamily\small,
  coltitle=white,
  title={#1},
  arc=2mm,
  boxrule=0.8pt,
  left=9pt, right=9pt, top=7pt, bottom=7pt
}

\usepackage{booktabs}
\usepackage{longtable}
\usepackage{array}
\usepackage{calc}

\usepackage{hyperref}
\hypersetup{
  colorlinks=true,
  linkcolor=accentblue,
  urlcolor=accentblue,
  citecolor=brandblue,
  pdfborder={0 0 0}
}

\title{\Huge\bfseries\sffamily\color{brandblue} T20 --- Java I/O, Streams, File, NIO.2 e Serializzazione \\[0.3em] \large\sffamily Byte Streams vs Character Streams, Bufferizzazione, Decorator Pattern I/O e Serializzazione}
\author{\textbf{Docente:} Prof. Michele Pasqua \quad---\quad \textbf{Appunti Revisione:} Wally}
\date{A.A. 2025/2026 -- Universit\`a degli Studi di Verona}

\begin{document}
\maketitle
\thispagestyle{fancy}
\vspace{-1.2em}
\noindent\textcolor{brandblue}{\rule{\textwidth}{0.8pt}}
\vspace{1.2em}

\begin{center}\rule{0.5\linewidth}{0.5pt}\end{center}

\subsection{1. Analisi Teorica Approfondita (Slide
T20)}\label{analisi-teorica-approfondita-slide-t20}

La lezione 20 conclude il percorso teorico del corso esplorando il
sottosistema di \textbf{Input/Output (I/O)} di Java
(\passthrough{\lstinline!java.io!} e
\passthrough{\lstinline!java.net!}), la persistenza su disco, la
comunicazione di rete, e la gestione di formati dati strutturati moderni
(\textbf{CSV} e \textbf{JSON}) tramite librerie terze gestite con Maven.

\begin{center}\rule{0.5\linewidth}{0.5pt}\end{center}

\subsubsection{1.1 L'Astrazione di I/O Stream (Slide
3-4)}\label{lastrazione-di-io-stream-slide-3-4}

In Java, qualsiasi trasferimento di dati da o verso l'esterno è
modellato tramite l'astrazione di \textbf{I/O Stream} (\emph{flusso
sequenziale unidirezionale di dati}): * \textbf{Differenza tassonomica
cruciale}: Gli \textbf{I/O Streams} (\passthrough{\lstinline!java.io!})
non vanno confusi con gli \textbf{Stream funzionali} di Java 8
(\passthrough{\lstinline!java.util.stream.Stream!}). Gli I/O Stream
trasportano byte o caratteri fisici verso periferiche o file; gli Stream
di Java 8 elaborano pipeline di trasformazione di oggetti in memoria. *
Un I/O Stream può essere agganciato a: * File su disco. * Flussi
standard di processo: \passthrough{\lstinline!System.in!} (standard
input), \passthrough{\lstinline!System.out!} (standard output),
\passthrough{\lstinline!System.err!} (standard error). * Connessioni di
rete (\emph{socket} TCP/IP o endpoint HTTP). * Buffer di memoria (array
di byte o stringhe).

\begin{center}\rule{0.5\linewidth}{0.5pt}\end{center}

\subsubsection{1.2 Il Dualismo dell'I/O in Java: Byte Streams vs
Character Streams (Slide
4-9)}\label{il-dualismo-dellio-in-java-byte-streams-vs-character-streams-slide-4-9}

L'architettura di \passthrough{\lstinline!java.io!} è rigorosamente
bipartita in due gerarchie parallele:

\needspace{7cm}
\begin{asciibox}
                      ┌─────────────────────────┐
                      │ java.io Stream Classes  │
                      └────────────┬────────────┘
                                   │
             ┌─────────────────────┴─────────────────────┐
             ▼                                           ▼
  ┌─────────────────────┐                     ┌─────────────────────┐
  │    Byte Streams     │                     │  Character Streams  │
  │  (8-bit dati grezzi)│                     │ (16-bit Unicode UTF)│
  └──────────┬──────────┘                     └──────────┬──────────┘
             │                                           │
  ┌──────────┴──────────┐                     ┌──────────┴──────────┐
  ▼                     ▼                     ▼                     ▼
┌─────────────┐ ┌──────────────┐       ┌─────────────┐ ┌─────────────┐
│ InputStream │ │ OutputStream │       │   Reader    │ │   Writer    │
└──────┬──────┘ └──────┬───────┘       └──────┬──────┘ └──────┬──────┘
       │               │                      │               │
┌──────┴──────┐ ┌──────┴───────┐       ┌──────┴──────┐ ┌──────┴──────┐
│FileInput-   │ │FileOutput-   │       │ FileReader  │ │ FileWriter  │
│Stream       │ │Stream        │       └──────┬──────┘ └──────┬──────┘
└──────┬──────┘ └──────┬───────┘              │               │
       │               │               ┌──────┴──────┐ ┌──────┴──────┐
┌──────┴──────┐ ┌──────┴───────┐       │Buffered-    │ │Buffered-    │
│Buffered-    │ │Buffered-     │       │Reader       │ │Writer       │
│InputStream  │ │OutputStream  │       └─────────────┘ └─────────────┘
└─────────────┘ └──────────────┘
\end{asciibox}

\paragraph{\texorpdfstring{1. Byte Streams (\texttt{InputStream} e
\texttt{OutputStream}, Slide
5-7)}{1. Byte Streams (InputStream e OutputStream, Slide 5-7)}}\label{byte-streams-inputstream-e-outputstream-slide-5-7}

\begin{itemize}
\tightlist
\item
  \textbf{Unità di dato}: Singolo byte grezzo (8 bit, intervallo
  \([0, 255]\) restituito come \passthrough{\lstinline!int!}, dove il
  valore \(-1\) segnala la fine dello stream --- \emph{End Of File,
  EOF}).
\item
  \textbf{Destinazione d'uso}: Immagini, file multimediali, bytecode
  compilato (\passthrough{\lstinline!.class!}), file compressi
  (\passthrough{\lstinline!.zip!}), comunicazioni binarie di rete.
\item
  \textbf{Metodi cardine}:

  \begin{itemize}
  \tightlist
  \item
    \passthrough{\lstinline!int read()!}: legge il prossimo byte (o
    \(-1\) a fine stream).
  \item
    \passthrough{\lstinline!int read(byte[] b)!}: riempie il buffer di
    byte.
  \item
    \passthrough{\lstinline!void write(int b)!}: scrive un byte.
  \item
    \passthrough{\lstinline!void write(byte[] b)!}: scrive un intero
    array di byte.
  \item
    \passthrough{\lstinline!void close()!}: rilascia il descrittore del
    file del sistema operativo.
  \end{itemize}
\item
  \textbf{Implementazioni notevoli}:

  \begin{itemize}
  \tightlist
  \item
    \passthrough{\lstinline!FileInputStream!} /
    \passthrough{\lstinline!FileOutputStream!}: lettura/scrittura
    diretta su disco.
  \item
    \passthrough{\lstinline!BufferedInputStream!} /
    \passthrough{\lstinline!BufferedOutputStream!}: aggiunge un buffer
    in memoria per minimizzare le costose chiamate di sistema del kernel
    del SO.
  \item
    \passthrough{\lstinline!DataInputStream!} /
    \passthrough{\lstinline!DataOutputStream!}: serializzazione di
    primitivi Java (\passthrough{\lstinline!readInt()!},
    \passthrough{\lstinline!writeDouble()!}).
  \end{itemize}
\end{itemize}

\paragraph{\texorpdfstring{2. Character Streams (\texttt{Reader} e
\texttt{Writer}, Slide
8-10)}{2. Character Streams (Reader e Writer, Slide 8-10)}}\label{character-streams-reader-e-writer-slide-8-10}

\begin{itemize}
\tightlist
\item
  \textbf{Unità di dato}: Caratteri Unicode (16 bit UTF-16).
\item
  \textbf{Destinazione d'uso}: Esclusivamente file di testo, file
  sorgente, documenti testuali.
\item
  \textbf{Le classi ponte (\emph{Bridge Adapters})}:

  \begin{itemize}
  \tightlist
  \item
    \passthrough{\lstinline!InputStreamReader!}: converte un
    \passthrough{\lstinline!InputStream!} di byte in un
    \passthrough{\lstinline!Reader!} di caratteri applicando una
    specifica codifica (es. UTF-8).
  \item
    \passthrough{\lstinline!OutputStreamWriter!}: converte caratteri in
    byte da inviare a un \passthrough{\lstinline!OutputStream!}.
  \end{itemize}
\item
  \textbf{Implementazioni notevoli}:

  \begin{itemize}
  \tightlist
  \item
    \passthrough{\lstinline!FileReader!} /
    \passthrough{\lstinline!FileWriter!}: lettura/scrittura di file di
    testo.
  \item
    \passthrough{\lstinline!BufferedReader!}: legge blocchi di testo
    bufferizzati e fornisce il comodissimo metodo
    \passthrough{\lstinline!String readLine()!} (restituisce un'intera
    riga o \passthrough{\lstinline!null!} a fine file).
  \item
    \passthrough{\lstinline!BufferedWriter!}: scrittura testuale con
    supporto al metodo \passthrough{\lstinline!newLine()!}.
  \end{itemize}
\end{itemize}

\begin{center}\rule{0.5\linewidth}{0.5pt}\end{center}

\subsubsection{\texorpdfstring{1.3 Gestione delle Risorse: Da
\texttt{try-finally} a \texttt{try-with-resources} (Slide 7,
10-11)}{1.3 Gestione delle Risorse: Da try-finally a try-with-resources (Slide 7, 10-11)}}\label{gestione-delle-risorse-da-try-finally-a-try-with-resources-slide-7-10-11}

\paragraph{Il Vecchio Approccio (Pre-Java 7, Slide
10):}\label{il-vecchio-approccio-pre-java-7-slide-10}

\begin{lstlisting}[language=Java]
BufferedReader br = null;
try {
    br = new BufferedReader(new FileReader("i.txt"));
    // lettura...
} catch (IOException e) {
    // gestione errore...
} finally {
    if (br != null) {
        try {
            br.close(); // Ulteriore try-catch obbligatorio perché close() lancia IOException!
        } catch (IOException e) { ... }
    }
}
\end{lstlisting}

Questo pattern era verbose, faticoso da manutenere e fonte continua di
\emph{resource leaks} (descrittori di file aperti non chiusi se si
verificavano eccezioni nel \passthrough{\lstinline!finally!}).

\paragraph{\texorpdfstring{L'Approccio Moderno:
\texttt{try-with-resources} (Java 7+, Slide
11)}{L'Approccio Moderno: try-with-resources (Java 7+, Slide 11)}}\label{lapproccio-moderno-try-with-resources-java-7-slide-11}

Tutte le classi che implementano l'interfaccia standard
\passthrough{\lstinline!java.lang.AutoCloseable!} possono essere
dichiarate all'interno delle parentesi tonde del blocco
\passthrough{\lstinline!try!}:

\begin{lstlisting}[language=Java]
try (BufferedReader br = new BufferedReader(new FileReader("i.txt"));
     BufferedWriter bw = new BufferedWriter(new FileWriter("o.txt"))) {
    String line;
    while ((line = br.readLine()) != null) {
        bw.write(line);
        bw.newLine();
    }
} catch (FileNotFoundException fnfe) {
    System.out.println("File not found...");
} catch (IOException ioe) {
    System.out.println("I/O Error...");
}
\end{lstlisting}

\begin{itemize}
\tightlist
\item
  \textbf{Garanzie della JVM}:

  \begin{enumerate}
  \def\labelenumi{\arabic{enumi}.}
  \tightlist
  \item
    Le risorse vengono chiuse \textbf{automaticamente} non appena il
    blocco \passthrough{\lstinline!try!} termina, sia in caso di
    completamento regolare sia in caso di eccezione o
    \passthrough{\lstinline!return!} anticipato.
  \item
    Vengono chiuse in \textbf{ordine inverso} rispetto alla loro
    dichiarazione (prima \passthrough{\lstinline!bw!}, poi
    \passthrough{\lstinline!br!}).
  \item
    Se sia il corpo del \passthrough{\lstinline!try!} che la chiamata a
    \passthrough{\lstinline!close()!} sollevano eccezioni, l'eccezione
    del corpo è quella principale propagata, mentre l'eccezione di
    chiusura viene salvata come \textbf{Suppressed Exception}
    (recuperabile con \passthrough{\lstinline!e.getSuppressed()!}).
  \end{enumerate}
\end{itemize}

\begin{center}\rule{0.5\linewidth}{0.5pt}\end{center}

\subsubsection{1.4 Gestione di File e Risorse di Rete (URL/URI) (Slide
13-17)}\label{gestione-di-file-e-risorse-di-rete-urluri-slide-13-17}

\begin{itemize}
\tightlist
\item
  \textbf{La classe \passthrough{\lstinline!java.io.File!} (Slide
  13-14)}: Rappresenta il percorso astratto di un file o di una
  directory. Non legge né scrive dati direttamente, ma permette di
  interrogare e modificare il filesystem:

  \begin{itemize}
  \tightlist
  \item
    \passthrough{\lstinline!exists()!},
    \passthrough{\lstinline!isFile()!},
    \passthrough{\lstinline!isDirectory()!},
    \passthrough{\lstinline!length()!},
    \passthrough{\lstinline!getAbsolutePath()!}.
  \item
    \passthrough{\lstinline!mkdir()!},
    \passthrough{\lstinline!delete()!},
    \passthrough{\lstinline!renameTo(File dest)!},
    \passthrough{\lstinline!listFiles()!}.
  \end{itemize}
\item
  \textbf{URL e URI (\passthrough{\lstinline!java.net!}, Slide 15-17)}:

  \begin{itemize}
  \item
    \passthrough{\lstinline!URI!} modella l'identificatore formale
    (\emph{RFC 2396}).
  \item
    \passthrough{\lstinline!URL!} modella la locazione fisica e il
    protocollo per accedere alla risorsa web.
  \item
    \emph{Avvertenza contemporanea (Slide 16)}: Il costruttore
    \passthrough{\lstinline!new URL("...")!} è \textbf{deprecato} da
    Java 20. La prassi moderna impone di creare un'istanza di
    \passthrough{\lstinline!URI!} e convertirla:

\begin{lstlisting}[language=Java]
URL url = new URI("https://info.cern.ch/index.html").toURL();
InputStream in = url.openStream(); // Apre una connessione HTTP e scarica i dati
\end{lstlisting}
  \end{itemize}
\end{itemize}

\begin{center}\rule{0.5\linewidth}{0.5pt}\end{center}

\subsubsection{1.5 Dati Strutturati: CSV e JSON (Slide
18-26)}\label{dati-strutturati-csv-e-json-slide-18-26}

Java non possiede parser nativi integrati nel runtime per file CSV e
JSON. Per questi formati si ricorre a librerie esterne integrate tramite
dipendenze Maven.

\paragraph{1. File CSV con OpenCSV (Slide
19-22)}\label{file-csv-con-opencsv-slide-19-22}

\begin{itemize}
\item
  Un file CSV (\emph{Comma-Separated Values}) memorizza tabelle in
  formato testuale, separando le colonne con virgole o punti e virgola.
\item
  \textbf{Dipendenza Maven}:
  \passthrough{\lstinline!com.opencsv:opencsv:5.12.0!}.
\item
  \textbf{Scrittura}: \passthrough{\lstinline!CSVWriter!} scrive array
  di stringhe \passthrough{\lstinline!String[]!} gestendo
  automaticamente apici ed escape:

\begin{lstlisting}[language=Java]
CSVWriter csvw = new CSVWriter(new FileWriter("data.csv"));
csvw.writeNext(new String[]{ "id", "name", "address" });
\end{lstlisting}
\item
  \textbf{Lettura}: \passthrough{\lstinline!CSVReader!} con
  \passthrough{\lstinline!readNext()!} riga per riga o
  \passthrough{\lstinline!readAll()!}.
\end{itemize}

\paragraph{2. File JSON con Google Gson (Slide
23-26)}\label{file-json-con-google-gson-slide-23-26}

\begin{itemize}
\item
  JSON (\emph{JavaScript Object Notation}) è lo standard dominante per
  lo scambio dati basato su mappe chiave-valore
  \passthrough{\lstinline!\{\}!} e liste ordinate
  \passthrough{\lstinline![]!}.
\item
  \textbf{Dipendenza Maven}:
  \passthrough{\lstinline!com.google.code.gson:gson:2.13.2!}.
\item
  \textbf{Serializzazione e Deserializzazione Automatica (Data
  Binding)}:

\begin{lstlisting}[language=Java]
Gson gson = new Gson();

// 1. Oggetto Java -> Stringa JSON (Serializzazione)
Date date = new Date(1, 1, 1970);
String json = gson.toJson(date); // Restituisce '{"day":1,"month":1,"year":1970}'

// 2. Stringa JSON -> Oggetto Java (Deserializzazione)
Date d = gson.fromJson(json, Date.class); // Ricostruisce l'istanza valorizzando i campi!
\end{lstlisting}

  Gson accede ai campi (anche \passthrough{\lstinline!private!} e
  \passthrough{\lstinline!final!}) tramite \textbf{Reflection}, senza
  richiedere getter/setter pubblici o costruttori senza argomenti!
\item
  \textbf{Pretty Printing}:
  \passthrough{\lstinline!new GsonBuilder().setPrettyPrinting().create()!}
  formatta il JSON con ritorni a capo e indentazione a 2 spazi.
\end{itemize}

\begin{center}\rule{0.5\linewidth}{0.5pt}\end{center}

\subsection{\texorpdfstring{2. Analisi Dettagliata del Codice
(\texttt{Lezione19/MavenDate})}{2. Analisi Dettagliata del Codice (Lezione19/MavenDate)}}\label{analisi-dettagliata-del-codice-lezione19mavendate}

In \passthrough{\lstinline!Lezione19/MavenDate!} il prof. Pasqua integra
tutti i concetti della Slide T20 all'interno della classe dimostrativa
\href{file:///home/wally/Documenti/GitHub/Programmazione-II/25-26/Teoria-e-Esercitazioni/Codice-20260902/Lezione19/MavenDate/src/main/java/it/oop/ui/MainIO.java}{MainIO.java}.

\subsubsection{\texorpdfstring{2.1 Configurazione delle Dipendenze in
\href{file:///home/wally/Documenti/GitHub/Programmazione-II/25-26/Teoria-e-Esercitazioni/Codice-20260902/Lezione19/MavenDate/pom.xml\#L68-L78}{pom.xml}}{2.1 Configurazione delle Dipendenze in pom.xml}}\label{configurazione-delle-dipendenze-in-pom.xml}

\begin{lstlisting}[language=XML]
        <!-- OpenCSV per parsing e generazione CSV (Slide 20) -->
        <dependency>
            <groupId>com.opencsv</groupId>
            <artifactId>opencsv</artifactId>
            <version>5.12.0</version>
        </dependency>
        <!-- Google Gson per serializzazione/deserializzazione JSON (Slide 24) -->
        <dependency>
            <groupId>com.google.code.gson</groupId>
            <artifactId>gson</artifactId>
            <version>2.13.2</version>
        </dependency>
\end{lstlisting}

\begin{center}\rule{0.5\linewidth}{0.5pt}\end{center}

\subsubsection{\texorpdfstring{2.2 Analisi Riga per Riga di
\href{file:///home/wally/Documenti/GitHub/Programmazione-II/25-26/Teoria-e-Esercitazioni/Codice-20260902/Lezione19/MavenDate/src/main/java/it/oop/ui/MainIO.java}{MainIO.java}}{2.2 Analisi Riga per Riga di MainIO.java}}\label{analisi-riga-per-riga-di-mainio.java}

\paragraph{\texorpdfstring{1. Lettura Binaria di Bytecode con
\texttt{FileInputStream} e \texttt{FileChannel} (Righe
18-41)}{1. Lettura Binaria di Bytecode con FileInputStream e FileChannel (Righe 18-41)}}\label{lettura-binaria-di-bytecode-con-fileinputstream-e-filechannel-righe-18-41}

\begin{lstlisting}[language=Java]
        FileInputStream fis = null;
        try {
            // Apertura dello stream di byte su un file binario compilato (.class)
            fis = new FileInputStream("src/main/resources/Date.class");
            FileChannel fc = fis.getChannel();
            int b;
            // Lettura byte a byte
            while ((b = fis.read()) != -1) {
                System.out.println((char) b);
            }
            // Rewind dello stream riportando la posizione a 0 tramite FileChannel
            fc.position(0);
            StringBuilder sb = new StringBuilder();
            // Lettura dell'intero contenuto in un colpo solo con readAllBytes() (Java 9+)
            for (byte bb : fis.readAllBytes()) {
                sb.append((char) bb);
            }
            System.out.println("-----");
            System.out.println(sb.toString());
        } catch (IOException ioe) {
            System.out.println(ioe.getMessage());
        } finally {
            // Chiusura classica pre-Java 7 con try-catch protetto
            try {
                if (fis != null) fis.close();
            } catch (IOException ioe) {
                System.out.println(ioe.getMessage());
            }
        }
\end{lstlisting}

\paragraph{\texorpdfstring{2. Scrittura di Testo con \texttt{FileWriter}
e \texttt{try-with-resources} (Righe
42-50)}{2. Scrittura di Testo con FileWriter e try-with-resources (Righe 42-50)}}\label{scrittura-di-testo-con-filewriter-e-try-with-resources-righe-42-50}

\begin{lstlisting}[language=Java]
        // Blocco try-with-resources: fw viene chiuso automaticamente
        try (FileWriter fw = new FileWriter("src/main/resources/file.txt")) {
            StringBuilder sb = new StringBuilder();
            // Genera stringhe di 'a' crescenti usando Stream di Java 8
            Stream.iterate("a", s -> s + "a")
                    .limit(15)
                    .forEach(s -> sb.append(s).append("\n"));
            fw.write(sb.toString());
        } catch (IOException ioe) {
            System.out.println(ioe.getMessage());
        }
\end{lstlisting}

\paragraph{3. Generazione di Tabella CSV con OpenCSV (Righe
52-65)}\label{generazione-di-tabella-csv-con-opencsv-righe-52-65}

\begin{lstlisting}[language=Java]
        try (CSVWriter csvw = new CSVWriter(new FileWriter("src/main/resources/data.csv"))) {
            String[] row = { "id", "name", "address" };
            csvw.writeNext(row); // Scrive l'intestazione delle colonne

            row[0] = "3";
            row[1] = "Paul";
            row[2] = "Strada le Grazie, 15";
            csvw.writeNext(row);

            row[0] = "6";
            row[1] = "Sam";
            row[2] = "Strada le Grazie, 18";
            csvw.writeNext(row);
        } catch (IOException ioe) {
            System.out.println(ioe.getMessage());
        }
\end{lstlisting}

\paragraph{4. Serializzazione e Deserializzazione con Gson (Righe
66-80)}\label{serializzazione-e-deserializzazione-con-gson-righe-66-80}

\begin{lstlisting}[language=Java]
        Gson gson = new Gson();
        Date date = new Date(1, 1, 1970);

        // Serializzazione dell'oggetto Date in formato JSON
        String json = gson.toJson(date);
        System.out.println(json); // Stampa: {"day":1,"month":1,"year":1970}

        // Deserializzazione da stringa JSON a nuova istanza Date
        Date date1 = gson.fromJson("{\"day\":1,\"month\":1,\"year\":1971}", Date.class);
        System.out.println(date1.toString()); // Stampa: y1971m1d1

        // Scrittura su file con Pretty Printing abilitato
        try (FileWriter fw = new FileWriter("src/main/resources/file.json")) {
            new GsonBuilder().setPrettyPrinting().create().toJson(date, fw);
        } catch (IOException ioe) {
            System.out.println(ioe.getMessage());
        }
\end{lstlisting}

\begin{center}\rule{0.5\linewidth}{0.5pt}\end{center}

\subsection{3. Risultato di Compilazione ed Esecuzione
Reale}\label{risultato-di-compilazione-ed-esecuzione-reale}

Comando eseguito nel progetto
\href{file:///home/wally/Documenti/GitHub/Programmazione-II/25-26/Teoria-e-Esercitazioni/Codice-20260902/Lezione19/MavenDate}{Lezione19/MavenDate}:

\begin{lstlisting}[language=bash]
mvn compile exec:java -Dexec.mainClass="it.oop.ui.MainIO"
\end{lstlisting}

Output ottenuto a video:

\begin{lstlisting}
[Dumping del bytecode del file Date.class con header CAFEBABE e constant pool]
-----
{"day":1,"month":1,"year":1970}
y1971m1d1
[INFO] BUILD SUCCESS
\end{lstlisting}

\subsubsection{Ispezione dei file generati su
disco:}\label{ispezione-dei-file-generati-su-disco}

\begin{enumerate}
\def\labelenumi{\arabic{enumi}.}
\item
  \passthrough{\lstinline!src/main/resources/file.txt!}:

\begin{lstlisting}
a
aa
aaa
...
aaaaaaaaaaaaaaa
\end{lstlisting}
\item
  \passthrough{\lstinline!src/main/resources/data.csv!}:

\begin{lstlisting}
"id","name","address"
"3","Paul","Strada le Grazie, 15"
"6","Sam","Strada le Grazie, 18"
\end{lstlisting}
\item
  \passthrough{\lstinline!src/main/resources/file.json!}:

\begin{lstlisting}
{
  "day": 1,
  "month": 1,
  "year": 1970
}
\end{lstlisting}
\end{enumerate}

\begin{center}\rule{0.5\linewidth}{0.5pt}\end{center}

\subsection{Traguardo Raggiunto: Revisione Integrale del Corso
Completata!}\label{traguardo-raggiunto-revisione-integrale-del-corso-completata}

Abbiamo completato con successo e senza tralasciare alcun dettaglio: *
Tutte le \textbf{20 slide teoriche} (da \passthrough{\lstinline!T01!} a
\passthrough{\lstinline!T20!}). * Tutti i \textbf{progetti e le lezioni
pratiche} (da \passthrough{\lstinline!Lezione04!} a
\passthrough{\lstinline!Lezione19!}). * Tutti i diagrammi, meme, pattern
architetturali, finezze di linguaggio, bug del docente, test Maven e
logiche di compilazione ed esecuzione.

Fammi sapere se desideri approfondire specifici temi per l'esame,
affrontare vecchi temi d'esame scritti/pratici, o dedicarti agli
esercizi del tuo workspace
(\passthrough{\lstinline!esercizi-wally-25-26!})!


\end{document}
